\section{Finalidad del manual}

Este manual explica cómo funciona una instalación dentro de Allplan desde el punto de vista del usuario.

Se ha tomado la instalación de Agua como ejemplo, pero el funcionamiento general es común al resto de instalaciones. Lo que cambia entre unas y otras son algunos elementos concretos, ciertos nombres de layer y algunos atributos, pero la forma de trabajar es prácticamente la misma.

El objetivo de este documento es que cualquier usuario pueda entender:

\begin{itemize}
\item cómo empezar una instalación;
\item cómo dibujarla;
\item cómo editarla;
\item cómo aplicar layers y atributos;
\item cómo insertar elementos definidos, macros y soportes;
\item cómo funcionan las copias o duplicados;
\item y qué comprobaciones conviene hacer antes de finalizar.
\end{itemize}

\section{A quién va dirigido}

Este manual está pensado para usuarios que necesiten modelar instalaciones en Allplan y entender el comportamiento general de la herramienta, sin entrar en detalles de programación ni en nombres internos del sistema.

\section{Cómo se complementa este manual}

Este documento se organiza en dos grandes bloques:

\begin{enumerate}
\item El funcionamiento global de las instalaciones.
\item Las particularidades de cada instalación.
\end{enumerate}

En la primera parte se explica la lógica de uso común a todos los sistemas.

En la segunda parte se recogen las diferencias específicas de cada instalación, por ejemplo:

\begin{itemize}
\item los tipos de instalación disponibles;
\item las distribuciones admitidas;
\item los diámetros disponibles;
\item los accesorios automáticos;
\item los elementos definidos propios;
\item las limitaciones o avisos especiales;
\item y cualquier diferencia importante respecto a otras instalaciones.
\end{itemize}

De este modo, el usuario puede entender primero el funcionamiento general y después consultar, dentro del mismo manual, la unidad concreta de la instalación con la que vaya a trabajar.

Cuando aparezca el marcador \pendingtag, significa que ese dato todavía debe validarse en Allplan antes de cerrar la versión definitiva del manual.

\section{Qué partes son comunes en todas las instalaciones}

Todas las instalaciones comparten la misma lógica de trabajo:

\begin{enumerate}
\item Se selecciona el tipo de instalación.
\item Se configuran los parámetros principales en la paleta.
\item Se dibuja o edita el recorrido.
\item Se revisa la previsualización.
\item Se añaden, si hace falta, elementos definidos, macros o soportes.
\item Se aplican layers y atributos.
\item Se finaliza la creación.
\end{enumerate}

Lo que cambia entre instalaciones suele ser:

\begin{itemize}
\item el tipo de elementos disponibles;
\item los diámetros;
\item las conexiones;
\item algunos layers;
\item y algunos atributos asociados a cada sistema.
\end{itemize}

\section{Vista general del trabajo con una instalación}

El trabajo habitual con una instalación se basa en tres ideas:

\begin{itemize}
\item dibujar un recorrido;
\item enriquecer ese recorrido con información;
\item y generar el resultado final.
\end{itemize}

En la práctica, esto significa que primero se crea el trazado principal y después se añaden o ajustan el resto de elementos necesarios.

\manualimagesized{CAPTURA-01.png}{Paleta principal completa de la instalación con sus apartados visibles}{0.52\textwidth}{0.62\textheight}

\section{Flujo de trabajo recomendado}

El orden más recomendable para trabajar es este:

\begin{enumerate}
\item Elegir el tipo de instalación.
\item Elegir el tipo de tubo o sistema.
\item Definir diámetro, distribución y resto de parámetros básicos.
\item Dibujar el recorrido principal.
\item Revisar uniones, codos y bifurcaciones.
\item Aplicar layers y atributos si es necesario.
\item Insertar elementos definidos, macros o soportes.
\item Comprobar que la información esté completa.
\item Pulsar el botón de finalizar.
\end{enumerate}

Seguir este orden ayuda a evitar correcciones posteriores innecesarias.

\section{Modos principales de trabajo}

La herramienta suele trabajar con tres modos principales.

\subsection{Modo creación}

Es el modo que se utiliza para dibujar un recorrido nuevo.

En este modo el usuario va definiendo puntos y la instalación crea el trazado entre ellos.

\subsection{Modo edición}

Es el modo que se utiliza para corregir o modificar un recorrido ya dibujado.

En este modo se pueden seleccionar tramos, borrar partes, cambiar diámetros o ajustar la geometría.

\subsection{Modo extender}

Es el modo que se utiliza para continuar un recorrido ya existente sin empezar desde cero.

\section{Qué debe configurar el usuario antes de dibujar}

Antes de empezar a dibujar conviene revisar los campos principales de la paleta.

\subsection{Tipo de instalación}

Es el campo donde el usuario elige el sistema con el que va a trabajar.

En Agua, por ejemplo, puede elegir entre varias opciones como:

\begin{itemize}
\item Polietilè;
\item Multicapa;
\item Armaflex.
\end{itemize}

Este campo define qué geometría y qué comportamiento tendrá la instalación.

\manualimage{CAPTURA-02.png}{Desplegable del tipo de instalación abierto con las opciones visibles}

\subsection{Diámetro}

El diámetro condiciona:

\begin{itemize}
\item el tamaño del elemento principal;
\item las piezas de conexión compatibles;
\item y parte de los atributos asociados.
\end{itemize}

Antes de dibujar, conviene comprobar que el diámetro seleccionado es el correcto.

\subsection{Distribución}

La distribución define cómo debe comportarse la instalación dentro del sistema.

En Agua es habitual trabajar con modos como:

\begin{itemize}
\item TD;
\item IS.
\end{itemize}

La distribución puede afectar a:

\begin{itemize}
\item la forma de los elementos;
\item los atributos aplicados;
\item y el comportamiento de ciertos accesorios.
\end{itemize}

\subsection{Tipo de agua u otras variantes}

Algunas instalaciones incluyen campos adicionales para distinguir variantes internas del sistema. Si la paleta muestra este dato, conviene definirlo antes de empezar.

\subsection{Ángulos y orientación}

Si la instalación trabaja con limitación angular o con orientación definida, estos parámetros deben comprobarse al inicio para evitar tener que rehacer el recorrido más adelante.

\section{Cómo dibujar una instalación}

El dibujo de la instalación se basa en una secuencia de clics que definen el recorrido.

\subsection{Inicio del recorrido}

El usuario hace clic en el primer punto del trazado.

\subsection{Continuación del recorrido}

Cada clic adicional añade un nuevo tramo.

La herramienta va construyendo la instalación entre esos puntos y mostrando una previsualización del resultado.

\subsection{Qué ve el usuario mientras dibuja}

Mientras se dibuja, la herramienta no solo muestra la línea del recorrido, sino también el comportamiento previsto de la instalación:

\begin{itemize}
\item tramos principales;
\item cambios de dirección;
\item uniones;
\item bifurcaciones.
\end{itemize}

\manualimage{CAPTURA-03.png}{Ejemplo de instalación en fase de dibujo con previsualización activa}

\section{Cómo se generan automáticamente los accesorios}

Una vez definido el recorrido, la herramienta interpreta la geometría y coloca automáticamente los accesorios necesarios.

\subsection{Codos}

Cuando el recorrido cambia de dirección, la herramienta genera el codo correspondiente.

\subsection{Manguitos o uniones}

Cuando dos tramos deben unirse de forma alineada, la herramienta genera la unión correspondiente.

\subsection{Bifurcaciones o tes}

Cuando un punto conecta tres direcciones, la herramienta genera la bifurcación adecuada.

\subsection{Regla general de prioridad}

Si un punto funciona como bifurcación, prevalece la bifurcación sobre otras soluciones.

Esto significa que el sistema no coloca varias piezas incompatibles en el mismo punto.

\manualtripleimage{CAPTURA-04A.png}{CAPTURA-04B.png}{CAPTURA-04C.png}{Ejemplo comparativo de un codo, una unión y una bifurcación}

\section{Cómo editar una instalación ya dibujada}

Una vez creado el recorrido, el usuario puede modificarlo.

\subsection{Seleccionar tramos}

Se pueden seleccionar:

\begin{itemize}
\item tramos individuales;
\item varios tramos a la vez;
\item o zonas concretas del recorrido.
\end{itemize}

\subsection{Borrar una sección}

El botón de borrar elimina el tramo o la parte seleccionada.

Después del borrado, la herramienta vuelve a analizar la conexión entre los tramos que quedan y reajusta automáticamente las piezas necesarias.

\subsection{Cambiar el diámetro}
\label{sec:cambiar-diametro}

El botón de modificar diámetro actualiza el diámetro del tramo o de la selección activa.

Cuando se modifica un único tramo, la herramienta muestra una ventana de confirmación con:

\begin{itemize}
\item la ruta;
\item el segmento;
\item el diámetro anterior;
\item y el nuevo diámetro.
\end{itemize}

El usuario debe confirmar el cambio antes de que la geometría se regenere.

\manualimage{CAPTURA-AGUA-04.png}{Confirmación de cambio de diámetro sobre un único tramo}

Cuando hay varios tramos seleccionados, la ventana de confirmación informa de:

\begin{itemize}
\item cuántos segmentos se van a modificar;
\item qué diámetro o diámetros tenían antes;
\item y qué nuevo diámetro se va a aplicar.
\end{itemize}

\manualimage{CAPTURA-AGUA-05.png}{Confirmación de cambio de diámetro sobre varios tramos seleccionados}

Después de aceptar, la instalación vuelve a reconstruir el resultado para adaptarlo al nuevo valor.

Este recálculo puede afectar no solo al tramo seleccionado, sino también a las uniones, codos, reducciones, conexiones o bifurcaciones relacionadas con ese tramo.

Después de cambiar un diámetro, conviene revisar en pantalla la zona modificada, especialmente cuando el cambio afecta a:

\begin{itemize}
\item dos tramos consecutivos con diámetros distintos;
\item un giro resuelto con codo;
\item una bifurcación;
\item o una conexión automática entre accesorios.
\end{itemize}

\section{Layers}

El usuario puede trabajar con los layers definidos por defecto o aplicar otros manualmente.

\subsection{Layers por defecto}

Cada instalación ya trae una configuración base de layers.

\subsection{Aplicar layers manualmente}

Desde la paleta, el usuario puede seleccionar un layer y aplicar ese layer a los elementos seleccionados.

Esto es útil cuando se necesita clasificar el modelo de una forma concreta antes de finalizar.

\manualimage{CAPTURA-05.png}{Ejemplo del apartado de layers en la paleta y su aplicación sobre la instalación}

\section{Atributos}

Además de la geometría, la instalación puede llevar información adicional en forma de atributos.

\subsection{Atributos automáticos}

Parte de los atributos se asignan automáticamente según el tipo de instalación, el elemento y la configuración elegida.

\subsection{Atributos aplicados por el usuario}

La paleta también permite aplicar atributos manualmente a una selección.

Esto es útil cuando el usuario necesita completar o corregir información antes de crear el resultado final.

\subsection{Atributo padre}

En determinadas instalaciones existe un atributo principal que sirve para organizar y relacionar elementos. Este dato es importante porque también puede influir en las copias o duplicados automáticos.

\section{Qué revisar antes de finalizar}

Antes de pulsar el botón de finalizar, se recomienda comprobar:

\begin{itemize}
\item que el recorrido es correcto;
\item que los diámetros son los adecuados;
\item que los layers necesarios están aplicados;
\item que los atributos importantes están informados;
\item y que los elementos definidos están bien colocados.
\end{itemize}

Si falta información, la herramienta puede mostrar un aviso antes de continuar.

\manualimage{CAPTURA-06.png}{Mensaje de aviso previo a finalizar cuando falta información}

\section{Rotación y orientación}

Algunos elementos admiten rotación y orientación.

Esto se aplica especialmente a:

\begin{itemize}
\item elementos definidos;
\item macros;
\item ciertos accesorios que requieren una posición concreta.
\end{itemize}

Cuando la paleta muestra campos de rotación u orientación, el usuario puede utilizarlos para ajustar la posición final del elemento.

\manualimage{CAPTURA-07.png}{Ejemplo de orientación o rotación aplicada a un elemento}

\section{Macros}

La herramienta permite colocar macros de librería dentro de la instalación.

\subsection{Para qué sirven}

Las macros permiten insertar objetos ya preparados, por ejemplo elementos de librería que deben colocarse en puntos concretos del recorrido.

\subsection{Qué debe hacer el usuario}

El flujo habitual es:

\begin{enumerate}
\item Elegir el tipo de macro.
\item Elegir el archivo correspondiente.
\item Elegir dónde se colocará.
\item Ajustar altura y rotación si es necesario.
\item Insertar la macro.
\end{enumerate}

\subsection{Dónde puede colocarse}

La macro puede situarse:

\begin{itemize}
\item al inicio del recorrido;
\item al final;
\item o en un punto libre.
\end{itemize}

\subsection{Altura de la macro}

La altura puede definirse directamente o en relación con un local seleccionado, según el caso.

\manualimagesized{CAPTURA-08.png}{Bloque de la paleta para insertar macros}{0.52\textwidth}{0.62\textheight}

\manualimagesized{CAPTURA-09.png}{Ejemplo de macro previsualizada dentro de la instalación}{0.62\textwidth}{0.55\textheight}

\section{Elementos definidos}

Además del recorrido principal, la herramienta permite añadir elementos definidos propios de la instalación.

En Agua, por ejemplo, pueden existir elementos como válvulas, salidas o piezas concretas del sistema.

\subsection{Cómo se insertan}

Hay dos formas habituales:

\begin{itemize}
\item insertarlos sobre un punto del recorrido;
\item o insertarlos como punto libre.
\end{itemize}

\subsection{Qué debe hacer el usuario}

El usuario debe:

\begin{enumerate}
\item elegir el tipo de elemento;
\item seleccionar el punto de inserción;
\item ajustar rotación o altura si hace falta;
\item confirmar la colocación.
\end{enumerate}

\manualimagesized{CAPTURA-10.png}{Bloque de la paleta para insertar Elementos Definidos}{0.52\textwidth}{0.62\textheight}

\manualimage{CAPTURA-11.png}{Inserción de un elemento definido como punto libre}

\section{Puntos no definidos}

Los puntos no definidos sirven para marcar la lógica del recorrido sin colocar todavía un elemento final.

Son útiles para describir:

\begin{itemize}
\item inicios;
\item finales;
\item intermedios;
\item bifurcaciones;
\item conexiones entre caminos.
\end{itemize}

\subsection{Para qué sirven}

Ayudan a organizar recorridos complejos y a definir la intención del trazado antes de convertirlo en una solución completa.

\subsection{Puntos comunes}

Si varios caminos coinciden en un punto o pasan muy cerca, la herramienta puede interpretar que existe una relación entre ellos.

Esto ayuda a construir una topología coherente.

\manualimage{CAPTURA-12.png}{Ejemplo de puntos no definidos en distintos colores}

\manualimage{CAPTURA-13.png}{Ejemplo de punto común entre dos caminos}

\section{Soportes}

La herramienta también permite trabajar con soportes.

\subsection{Qué hace el usuario}

El flujo habitual es:

\begin{enumerate}
\item configurar el soporte en la paleta;
\item elegir el tipo de soporte y revisar sus parámetros;
\item pulsar el botón para insertarlo;
\item definir su posición con los puntos necesarios;
\item revisar la previsualización;
\item acumularlo o crearlo.
\end{enumerate}

Para que el soporte quede bien definido, conviene completar antes de insertarlo los campos de tipo, superficie, cotas y ángulo.

\subsection{Tipos de soporte}

De forma general, la herramienta permite trabajar con familias de soporte como:

\begin{itemize}
\item \texttt{Zeta}
\item \texttt{Omega}
\item \texttt{Cinta}
\end{itemize}

Cada una responde a una solución física distinta. Por eso, antes de colocarlo, el usuario debe elegir el tipo que corresponda al sistema y al montaje que quiere representar.

En función de la instalación, pueden aparecer variantes o subtipos asociados a cada familia.

\subsection{Subtipo de instalación}

Además del tipo de soporte, la paleta puede mostrar el campo \texttt{Subtipo instalación}.

Este campo ayuda a identificar con qué familia de instalación se relaciona el soporte. Es útil para mantener coherencia entre el soporte que se inserta y la instalación sobre la que se está trabajando.

\manualimagesized{CAPTURA-14B.png}{Tipo de soporte, subtipo instalación y superficie en la paleta de soportes}{0.60\textwidth}{0.55\textheight}

\subsection{Superficie}

El campo \texttt{Superficie} permite definir cómo debe considerarse el soporte:

\begin{itemize}
\item \texttt{Liso}
\item \texttt{Perforado}
\end{itemize}

Esta elección forma parte de la definición final del soporte y de la información que queda asociada a él. Por eso conviene revisarla antes de insertarlo.

No todos los tipos de soporte utilizan esta distinción exactamente de la misma manera, pero para el usuario la regla práctica es sencilla: debe escoger la superficie que corresponda al soporte real que quiere representar.

\subsection{Cota A y Cota B}

Estas dos medidas son básicas para definir la geometría del soporte:

\begin{itemize}
\item \texttt{Cota A}: indica la altura del soporte o su desarrollo vertical.
\item \texttt{Cota B}: indica el largo del soporte o su desarrollo horizontal.
\end{itemize}

En soportes con parte vertical, \texttt{Cota A} tiene un papel especialmente importante. En otros tipos, como ciertas cintas, no siempre interviene de la misma forma.

\subsection{Ángulo de inclinación}

El campo \texttt{Ángulo de inclinación} permite girar o inclinar el soporte respecto a su posición base.

Es útil cuando el soporte no debe quedar en una orientación neutra y necesita un ajuste adicional para adaptarse al montaje real.

Lo recomendable es definir primero la posición del soporte y usar después el ángulo como ajuste fino.

\manualimage{CAPTURA-14C.png}{Ejemplo de Cota A, Cota B y ángulo de inclinación en la paleta}

\subsection{Modos de edición}

Una vez que hay soportes acumulados, la herramienta permite trabajar con distintos modos de edición:

\begin{itemize}
\item \texttt{Desactivado}: la herramienta no entra en edición de soportes y los clics siguen el flujo normal de trabajo.
\item \texttt{Edición}: permite seleccionar soportes para revisarlos, borrarlos o aplicarles atributos.
\item \texttt{Edición Mover}: permite recolocar un soporte ya insertado o acumulado.
\end{itemize}

Estos modos son útiles cuando el recorrido principal ya está dibujado y solo falta ajustar la posición o la información de los soportes.

\manualimage{CAPTURA-14D.png}{Modos de edición de soportes en la paleta}

\subsection{Atributos de soporte}

La paleta también incluye un campo para aplicar el \texttt{Atributo de soporte}.

Este atributo sirve para identificar, clasificar o completar la información de los soportes antes de la creación final.

Si el usuario tiene soportes seleccionados, el atributo se aplica a esa selección. Si no hay una selección activa, la aplicación afecta a los soportes acumulados dentro de la operación en curso.

Esto resulta especialmente útil cuando se quiere dejar todos los soportes correctamente preparados antes de finalizar la instalación.

\manualimage{CAPTURA-14E.png}{Campo de atributo de soporte y aplicación sobre soportes seleccionados o acumulados}

\subsection{Cuándo conviene colocarlos}

Normalmente es más cómodo colocar los soportes cuando el recorrido principal ya está bastante definido.

También conviene tener en cuenta que, antes de acumular o crear un soporte, primero debe haberse posicionado correctamente. Si todavía no se ha definido su colocación, la herramienta avisa y no lo añade.

\manualimage{CAPTURA-14.png}{Bloque de soportes en la paleta y preview de un soporte}

\section{Cómo funcionan los duplicados o copias}

Este punto es importante porque puede generar dudas.

\subsection{Duplicados no deseados}

La herramienta evita repetir elementos iguales cuando no corresponde.

Esto ayuda a que no aparezcan piezas duplicadas por error en el resultado final.

\subsection{Copias intencionadas}

En algunos casos, la instalación puede generar copias en otros archivos de dibujo según la información asociada al elemento.

Estas copias no son un error. Forman parte del funcionamiento previsto de la herramienta.

De forma general, el comportamiento es este:

\begin{enumerate}
\item el elemento se crea en el archivo principal en el que el usuario está trabajando;
\item si ese elemento tiene informado un atributo padre, la herramienta puede buscar un archivo de dibujo relacionado con esa referencia;
\item si el archivo de destino está cargado y coincide con esa referencia, la herramienta crea allí una copia relacionada;
\item esa copia queda vinculada a la actualización posterior de la instalación.
\end{enumerate}

\subsection{Qué debe entender el usuario}

El usuario debe saber que:

\begin{itemize}
\item puede existir un elemento en el archivo principal;
\item y puede existir una copia relacionada en otro archivo de dibujo.
\end{itemize}

Para que esa copia pueda producirse, normalmente deben cumplirse estas condiciones:

\begin{itemize}
\item el atributo padre debe estar informado;
\item el archivo de dibujo de destino debe estar cargado;
\item y el nombre del archivo de destino debe coincidir con la referencia que utiliza la instalación.
\end{itemize}

Si el archivo de destino no está cargado o no coincide con esa referencia, la copia no se genera.

\subsection{Qué pasa cuando la instalación se vuelve a editar}

Cuando la instalación se actualiza, esas copias también se actualizan para que no queden versiones antiguas.

Antes de crear las nuevas copias, la herramienta localiza las anteriores, las elimina de los archivos correspondientes y vuelve a generar las copias actualizadas.

\manualtripleimagelabeled{CAPTURA-15A.png}{Archivo original en \texttt{TEST}}{CAPTURA-15B.png}{Copias relacionadas visibles en \texttt{TD01}}{CAPTURA-15C.png}{Copias relacionadas visibles en \texttt{IS01}}{Ejemplo del archivo original TEST y de las copias relacionadas generadas en TD01 e IS01}

\section{Qué ocurre al pulsar el botón de finalizar}

Cuando el usuario pulsa el botón de finalizar, la herramienta:

\begin{enumerate}
\item revisa la información disponible;
\item reconstruye los elementos necesarios;
\item aplica atributos y layers;
\item incorpora elementos definidos, macros o soportes;
\item crea el resultado definitivo;
\item y actualiza las copias relacionadas si las hubiera.
\end{enumerate}

En ese momento la instalación pasa de estar en fase de preparación a quedar creada como resultado final en el documento.

\section{Recomendaciones prácticas de uso}

\begin{itemize}
\item Configurar primero el sistema antes de empezar a dibujar.
\item No cambiar de criterio de diámetro continuamente durante el dibujo si no es necesario.
\item Revisar la distribución antes de finalizar.
\item Aplicar layers y atributos cuando el recorrido ya esté claro.
\item Insertar macros y elementos definidos cuando la base principal esté estable.
\item Revisar los avisos antes de aceptar la creación final.
\end{itemize}

\section{Resumen final}

La herramienta de instalaciones está pensada para que el usuario dibuje un recorrido, lo complete con la información necesaria y genere un resultado final coherente dentro de Allplan.

Aunque cada instalación tenga sus particularidades, el funcionamiento general es siempre el mismo:

\begin{itemize}
\item configurar;
\item dibujar;
\item revisar;
\item completar;
\item y finalizar.
\end{itemize}

Por eso, entendiendo bien el ejemplo de Agua, se entiende también el comportamiento general del resto de instalaciones.

